\documentclass[10pt]{article}
\usepackage[margin=2.2cm]{geometry}
\usepackage{graphicx}
\usepackage{amsmath}
\usepackage{booktabs}
\usepackage[hidelinks]{hyperref}
\setlength{\parskip}{4pt}\setlength{\parindent}{0pt}

\title{\vspace{-1.2cm}Neuron mechanisms II: what $d8$-dominant neurons compute\\
\large \texttt{wd3\_k0\_N2048}, step 500k --- post-nonlinearity units, uniform evaluation data}
\author{}\date{September 2026}

\begin{document}\maketitle\vspace{-1.4cm}

\section*{Setup}
We use the same protocol as in the $d4$ report: $a=\phi(z)$, $\phi=\mathrm{relu}(\tanh)$, $z=(Vh)_n$,
analyzed on all $2^{16}$ uniform inputs. A neuron is called \emph{$d8$-dominant} when the largest
Walsh coefficient among the tracked tree modes is either $\chi_{d8_0}$ (bits 0--7) or
$\chi_{d8_1}$ (bits 8--15). There are 2047 such neurons (307/452/619/669 in blocks 0--3).
For each neuron, we decompose the pre-activation $z$ in two complementary ways: over the full
256-mode Walsh basis of the neuron's own 8-bit half-cube, and over the three tree-level modes
$y_0=\chi_{d8_0}$, $y_1=\chi_{d8_1}$, and $\chi_{d16}=y_0y_1$. Looking at where the variance
of $z$ is concentrated cleanly separates the neurons into three mechanisms.

\section*{Three mechanisms, ordered by depth}
\textbf{(1) Blocks 0--1: 8-input majority gates on raw bits} (951 neurons; $100\%$ of
block~0, $94\%$ of block~1). Their pre-activation is well approximated by
$z\approx\gamma\sum_{k=1}^{8}s_kx_k$, with $\gamma\approx0.13$, and at least $99\%$ of
$\mathrm{Var}(z)$ lies in the eight degree-1 singleton modes. This is the same mechanism as
for the $d4$ neurons, but now applied to eight raw bits instead of four. The measured response
as a function of the signed input sum, $E[a\,|\,s]$, matches the ideal
$\mathrm{relu}(\tanh(\gamma s))$ gate essentially exactly (Fig.~\ref{fig:mech}, left).
Because such a gate has Walsh power spread over many modes of the 8-bit cube, only a small
fraction of its variance lies in the tracked binary-tree modes: the ideal coverage is $0.042$,
and the measured median is $0.044$. This continues the same trend seen at lower degree:
$d2$ corner units have coverage $1/3$, $d4$ gates have coverage $\approx0.15$, and $d8$ gates
have coverage $\approx0.044$. Figure~\ref{fig:gatespec} shows the full 255-mode Walsh
spectrum such a gate produces, in direct analogy to the $d4$ spectrum figure: the eight
singletons carry $\approx9\%$ each, all 28 pairs are equal at $\approx0.8\%$ (the four tree
pairs are in no way special), odd degrees are absent, and the $d8$ mode itself is a single
$\approx1.1\%$ spike.

\textbf{(2) Blocks 2--3: $8{+}8\!\to\!16$ corner units} (189 neurons). Their pre-activation is
a nearly perfectly balanced sum of the two already-formed $d8$ parities,
$z\approx\gamma(\pm y_0\pm y_1)$, with the variance split $49.7\%/49.8\%$ between the two
inputs and essentially no $d16$ component already present in $z$. After the nonlinearity,
these neurons become the same two-input corner detectors discussed in the $d2$ case:
$a=\frac{t}{4}(1\pm y_0\pm y_1\pm y_0y_1)$. In other words, the neuron responds strongly to
one particular sign combination of $(y_0,y_1)$ (for example, $0.42$ on $({+},{+})$ and
$\leq0.005$ on the other three corners; Fig.~\ref{fig:corner}). Its variance then splits into
\textbf{approximately equal thirds over $\{\chi_{d8_0},\chi_{d8_1},\chi_{d16}\}$}
($0.332/0.333/0.319$ for the exemplar). \emph{These neurons are therefore the point where the
degree-16 target is actually created.} They are labeled ``$d8$-dominant'' only because the
three coefficients are almost tied and small sampling differences make one of the $d8$ modes
slightly larger. Since all three modes belong to the tracked tree, their total tree coverage is
$\approx1.0$.

\textbf{(3) Blocks 2--3: amplifiers} (296 neurons). For these neurons, at least $90\%$ of the
pre-activation variance already lies in the neuron's own $d8$ feature assembled upstream, so
the neuron mainly strengthens or passes along an existing $d8$ representation rather than
constructing it from scratch. Their median tree coverage is $0.98$. The remaining ``other''
class (600 neurons, median coverage $0.93$) forms a continuum between the clean amplifier and
corner mechanisms: these neurons still carry substantial $d8$/$d16$ structure, but with a
larger contribution from additional modes, plus a small noisy tail. In fact, much of the
``other'' class is \emph{effectively dead}: 487 of the 2047 $d8$-dominant neurons have total
activation Walsh power $\mathrm{Var}(a)<10^{-4}$ (28\%/45\% of blocks 2--3), and 380 of them
are ``other''. Figure~\ref{fig:dead} repeats the composition of Fig.~\ref{fig:mech} (right)
with this taken into account: once neurons are weighted by their amplitude $\mathrm{Var}(a)$,
the picture sharpens --- blocks 0--2 are dominated by raw-bit gates, and block 3 by the
$8{+}8$ corner units ($64\%$), whose class total is carried largely by a few dozen strong
units.

\begin{figure}[h]\centering
\includegraphics[width=0.78\textwidth]{fig3_d8_mechanisms_nopanel.pdf}
\caption{Left: a block-0 $d8$ gate's measured response as a function of signed input sum,
compared with the ideal $\mathrm{relu}(\tanh(\gamma s))$ curve (dashed; the two are essentially
indistinguishable). Right: mechanism composition by depth --- raw-bit gates create $d8$
early, while corner units and amplifiers become dominant in blocks 2--3.}
\label{fig:mech}
\end{figure}

\begin{figure}[h]\centering
\includegraphics[width=\textwidth]{fig6_d8_gate_spectrum.pdf}
\caption{Activation Walsh spectrum of an 8-input raw-bit gate (block 0, neuron 1) over all
255 modes of its 8-bit cube, measured (bars) vs.\ the parameter-free ideal gate (ticks),
grouped by degree (the empty odd degrees are omitted). Colors mark the tree-tracked modes; the tree pairs and quads sit at
exactly the same height as their off-tree siblings, which is why tree coverage is only
$\approx0.044$.}
\label{fig:gatespec}
\end{figure}

\begin{figure}[h]\centering
\includegraphics[width=0.72\textwidth]{fig4_d8_corner.pdf}
\caption{An $8{+}8\!\to\!16$ corner unit (block 3, \#333). Left: the neuron fires almost only
for one of the four possible sign combinations of the two $d8$ parities. Right: its variance
splits into approximately equal thirds across the two $d8$ modes and the $d16$ mode --- the
$d16$ component is created by this neuron.}
\label{fig:corner}
\end{figure}

\begin{figure}[h]\centering
\includegraphics[width=\textwidth]{fig7_d8_composition_dead.pdf}
\caption{The composition of Fig.~\ref{fig:mech} (right), accounting for effectively-dead
neurons ($\mathrm{Var}(a)<10^{-4}$). Left: amplitude distribution by class --- 487 of the
2047 $d8$-dominant neurons are dead, mostly ``other''. Middle: the same counts with the dead
split out ($28\%$/$45\%$ of blocks 2--3). Right: composition weighted by $\mathrm{Var}(a)$:
the block-3 amplitude is carried mainly by the $8{+}8$ corner units ($64\%$, dominated by a
heavy tail --- the top 10 corners hold half the class total), while the dead sliver carries
$\leq1.3\%$ everywhere.}
\label{fig:dead}
\end{figure}

\begin{table}[h]\centering\small
\begin{tabular}{lccccc}
\toprule
class & input ($z$) & blocks & $n$ & tree coverage (median) \\
\midrule
majority gate & $\gamma\sum_8 s_kx_k$ (raw bits) & 0--1 (some 2) & 951 & 0.044 \\
$8{+}8$ corner & $\gamma(\pm y_0 \pm y_1)$ & 2--3 & 189 & 1.00 \\
amplifier & own $\chi_{d8}$ feature & 2--3 & 296 & 0.98 \\
other (mixtures) & partial $d8$/$d16$ + bulk & 2--3 & 600 & 0.93 \\
\bottomrule
\end{tabular}
\end{table}

\textbf{Takeaway.} The network does \emph{not} appear to build $d8$ by composing already-formed
$d4$ features. Instead, blocks 0--1 create $d8$ directly from raw degree-1 inputs using wide
8-input majority gates. This is the same basic mechanism seen for $d4$, but because the gate
acts on more inputs, an even smaller fraction of its Walsh spectrum lies in the tracked tree
($\approx0.04$). The only clear compositional step in this network occurs at the final level:
dedicated corner units take the two assembled $d8$ parities as inputs and create
$\chi_{d16}$ in blocks 2--3. This agrees with the layer-level decode, where the
$8{+}8\,|\,16$ partition becomes dominant at exactly those blocks. Thus the same basic
gate/corner mechanism appears across degrees; what changes is which features feed into it
(raw bits at lower levels versus already-assembled parities at the final level) and how much
of the resulting Walsh spectrum falls inside the binary-tree modes we track.

\vfill\footnotesize
Artifacts: \texttt{gdrive:parity/analysis\_output/neuron\_mechanisms/wd3\_k0\_N2048/step\_00500000/}
(\texttt{d8\_classes.csv}, \texttt{d8\_per\_neuron.csv}, spectra, figures, this report;
companion $d4$ and $d2$ probes in the same folder).
\end{document}
